\section{Introducción}

\subsection{Planteamiento del Problema}

En los establecimientos del sector restaurantero, la gesti\'on de pedidos y el proceso
de pago representan una parte fundamental de la experiencia del cliente y del control
operativo del negocio. Sin embargo, en muchos restaurantes estos procesos a\'un dependen
de m\'etodos tradicionales que implican la intervenci\'on constante del personal, registros
manuales o sistemas parcialmente digitalizados \cite{toast_success_2023, nra_state_2023}. Esto genera inconvenientes como errores
en los pedidos, retrasos en la comunicaci\'on con cocina y dificultades al momento de dividir
la cuenta entre varios comensales.

Uno de los problemas m\'as frecuentes ocurre cuando un grupo de personas comparte una
mesa y desea pagar \'unicamente por los productos que consumi\'o. En estos casos, el personal
del restaurante debe identificar manualmente los consumos individuales, calcular proporciones
o dividir la cuenta total, lo cual resulta tedioso, propenso a errores y genera tiempos de espera
innecesarios \cite{lee_digital_2024, tan_tabletop_2020}.

En M\'exico, el sector de servicios de preparaci\'on de alimentos y bebidas concentra m\'as de
700{,}000 unidades econ\'omicas de acuerdo con el Directorio Estad\'istico Nacional de Unidades
Econ\'omicas (DENUE) del Instituto Nacional de Estad\'istica y Geograf\'ia (INEGI) \cite{inegi_denue_2024}.
Asimismo, la C\'amara Nacional de la Industria de Restaurantes y Alimentos Condimentados (CANIRAC)
proyect\'o un crecimiento del 4\% para 2024 \cite{inegi_canirac_2023} y de entre 5.5\% y 6\% para 2025 \cite{ramirez_canirac_2025},
impulsado por la recuperaci\'on del consumo y la apertura de nuevas unidades.
A pesar de esta expansi\'on, la brecha tecnol\'ogica entre los sistemas existentes y las expectativas
de los comensales sigue siendo significativa: pocas soluciones integran en una misma plataforma
la interacci\'on directa del comensal con el men\'u, la gesti\'on de \'ordenes y la divisi\'on
aut\'onoma de pagos.

En los \'ultimos a\~nos han surgido diferentes sistemas tecnol\'ogicos orientados a la gesti\'on
de restaurantes, tales como sistemas de punto de venta (POS), aplicaciones de pedidos digitales
o plataformas de administraci\'on de inventarios. No obstante, muchas de estas soluciones se
centran principalmente en la operaci\'on interna del restaurante y no integran completamente
la participaci\'on directa del cliente en el proceso de selecci\'on de productos, gesti\'on de
pedidos y divisi\'on de pagos dentro de una misma plataforma \cite{katagri_smart_2025}.

Ante esta problem\'atica, el presente trabajo propone el desarrollo de un sistema omnicanal
basado en una aplicaci\'on web progresiva (PWA) que permita la creaci\'on de mesas virtuales
mediante c\'odigos QR, a trav\'es de las cuales los comensales puedan acceder al men\'u digital
del restaurante, seleccionar los productos que desean consumir y gestionar el pago de su consumo
de manera individual o colaborativa.

\subsection{Objetivos}

\subsubsection{Objetivo General}

Diseñar y desarrollar una plataforma web progresiva (PWA) de gestión de pedidos,
liquidación de cuentas y procesamiento analítico para cadenas restauranteras,
que permita a los comensales realizar pedidos individuales y compartidos desde
su propio dispositivo móvil mediante una mesa virtual protegida por contraseña,
mientras el personal operativo (meseros, chefs y gerentes) administra la operación
en tiempo real e integra un pipeline de procesamiento batch con Apache Spark
para la generación de métricas de negocio y pronósticos de demanda, desplegado
sobre una infraestructura en la nube basada en Vercel y Supabase.

\subsubsection{Objetivos Específicos}

Los siguientes objetivos específicos describen los módulos y componentes
funcionales que deben desarrollarse e integrarse para alcanzar el objetivo general.
Cada uno incluye la evidencia de verificación que confirma su cumplimiento.

\noindent\textbf{Módulo 1 — Infraestructura y Estructura Organizacional}

\noindent\textbf{OE-01 Infraestructura.}
Configurar el ecosistema en la nube integrado por Vercel para el despliegue de
la PWA, Supabase para la persistencia transaccional y un servidor privado virtual
(VPS) dedicado para el procesamiento distribuido con Apache Spark. Definir el
schema inicial en Prisma ORM y la estructura de almacenamiento para la arquitectura
Medallion (capas Bronze, Silver, Gold) dentro del entorno del VPS. Establecer el
pipeline de CI/CD para la automatización de despliegues y migraciones.
\textit{Evidencia: Entornos de producción y preview operativos en Vercel. Instancias
de Supabase activas con políticas RLS. VPS configurado con el runtime de Spark
y estructura de directorios Medallion lista para ingesta de datos.}

\noindent\textbf{OE-02 Estructura Organizacional.}
Desarrollar los módulos de creación, edición y visualización de los tres niveles
organizacionales: Cadena restaurantera (nivel~1), Zona (nivel~2) y Sucursal
(nivel~3). Cada nivel debe contener las entidades que le corresponden y garantizar
que toda entidad inferior esté siempre ligada a su entidad superior. El acceso a
cada nivel debe estar restringido conforme al rol y nivel organizacional del usuario
autenticado.
\textit{Evidencia: CRUD completo de cadena, zonas y sucursales. Validación de
integridad referencial en base de datos. Tests de acceso: un gerente de zona no
puede ver ni modificar sucursales de otra zona.}

\noindent\textbf{Módulo 2 — Autenticación, Usuarios y Control de Acceso}

\noindent\textbf{OE-03 Gobernanza de Datos: Autenticación y Control de Acceso.}
Establecer los mecanismos de gobernanza de datos mediante la implementación de un
sistema de autenticación robusto para usuarios colaboradores (Admin, Gerente de
Zona, Gerente de Sucursal, Mesero, Chef) utilizando Supabase Auth. Definir
políticas estrictas de Row Level Security (RLS) que garanticen la integridad,
privacidad y el principio de mínimo privilegio (\textit{PoLP}), asegurando que el
acceso a la información esté segmentado según el nivel organizacional y el rol
asignado.
\textit{Evidencia: Matriz de trazabilidad de permisos y roles validada. Políticas
RLS activas que impiden fugas de datos entre niveles organizacionales. Flujo de
primer inicio de sesión con cambio obligatorio de contraseña funcional.}

\noindent\textbf{OE-04 Gestión de Usuarios.}
Implementar el módulo de alta de usuarios donde el usuario habilitado ingresa el
correo del nuevo colaborador, el sistema genera una contraseña temporal, y se asigna
el rol según las reglas de negocio (un usuario no puede asignar un rol superior al
propio). El módulo debe filtrar los niveles y sucursales disponibles para la
asignación según el alcance del usuario creador.
\textit{Evidencia: Admin da de alta a un Gerente de Zona (exitoso). Gerente de
Sucursal intenta dar de alta a otro Gerente de Zona (rechazado). Nuevo usuario
realiza primer login y cambia contraseña correctamente.}

\noindent\textbf{Módulo 3 — Catálogo: Categorías y Productos}

\noindent\textbf{OE-05 Catálogo.}
Desarrollar los formularios y lógica de negocio para la gestión de catálogos en dos
niveles: Plantillas canónicas a nivel de cadena (\textit{Templates}) para
estandarización, e instancias locales por sucursal. El catálogo debe incluir
categorías, productos con variantes de precio, estaciones de preparación para el
KDS y grupos de modificadores (\textit{Extras}, \textit{Términos}) con reglas de
selección obligatoria u opcional.
\textit{Evidencia: CRUD completo de templates y catálogo local. Modificadores
asignados y validados. Productos enrutados correctamente a sus estaciones. Herencia
de impuestos por categoría verificada.}

\noindent\textbf{Módulo 4 — Gestión de Mesas Físicas y Virtuales}

\noindent\textbf{OE-06 Mesas Físicas.}
Implementar la pantalla de gestión de mesas físicas accesible para meseros y
gerentes, que permita definir la capacidad, número y ubicación de cada mesa en un
plano interactivo (\textit{Floor Map}) mediante coordenadas $X, Y$ y formas
geométricas. El sistema debe soportar la unión de múltiples mesas físicas en una
sola sesión operativa para grupos grandes.
\textit{Evidencia: Mapa de mesas funcional con representación visual de estatus.
Unión y separación de mesas verificada en la base de datos (\texttt{DiningSessionTable}).
Cambios de posición persistidos correctamente.}

\noindent\textbf{OE-07 Mesa Virtual.}
Desarrollar el flujo de inicialización de mesa virtual por parte del mesero/gerente
(selección de mesas físicas libres y generación de código QR con JWT de un solo
uso), y el flujo de activación por parte del anfitrión (escaneo del QR, creación
de contraseña y acceso al menú del restaurante). La mesa virtual debe pasar al
estado \textit{Por liquidar} al ser activada.
\textit{Evidencia: QR generado con JWT de un solo uso y TTL de 30 minutos.
Anfitrión activa mesa exitosamente. Comensal se une a la mesa con ID y contraseña correctos.}

\noindent\textbf{Módulo 5 — Órdenes y Cocina}

\noindent\textbf{OE-08 Órdenes.}
Implementar la interfaz de menú para comensales donde pueden explorar el catálogo,
agregar productos con cantidades a su canasta y confirmar la orden como individual
(cargada a su cuenta personal) o compartida (cargada a la mesa grupal). Cada
confirmación debe presentar un diálogo de verificación antes de enviar la orden a
cocina. Cada ítem de orden debe generar un \textit{snapshot} inmutable (nombre,
precio, variantes y modificadores) para garantizar la integridad histórica frente
a cambios futuros en el catálogo.
\textit{Evidencia: Comensal genera orden individual y aparece en su cuenta personal.
Comensal genera orden compartida y aparece en la cuenta grupal con su nombre
asociado. Orden rechazada si la mesa está en estado \textit{Liquidada}.}

\noindent\textbf{OE-09 Tablero Cocina.}
Desarrollar la vista de la Cocina donde se muestran todas las órdenes activas de la
sucursal con su estatus (Pendiente, En preparación, En reparto, Entregada) y la
posibilidad de actualizarlo. Las nuevas órdenes deben notificarse en tiempo real
vía Supabase Realtime con indicador visual y/o sonoro.
\textit{Evidencia: Nueva orden aparece en el tablero de la Cocina en menos de 2 segundos
sin recargar la página. Chef cambia estatus de orden y el cambio se refleja en la
vista del comensal en tiempo real. Órdenes de otras sucursales no visibles.}

\noindent\textbf{Módulo 6 — Liquidación y Pagos}

\noindent\textbf{OE-10 Liquidación.}
Implementar la pantalla de Órdenes y Liquidaciones donde el comensal visualiza el
desglose de cuentas individuales y el total de la mesa compartida. El comensal puede
seleccionar qué cuentas individuales liquidar, qué monto aportar a la mesa
compartida, o liquidar la mesa completa. El sistema debe validar que no se procesen
pagos concurrentes duplicados mediante control de concurrencia
(\texttt{SELECT FOR UPDATE} en PostgreSQL).
\textit{Evidencia: Liquidación de cuenta individual exitosa. Intento de doble pago
simultáneo: solo uno es procesado, el segundo recibe error descriptivo. Monto de
mesa compartida no superable. Todos los comensales se muestran como
\textit{Liquidados} al completarse sus pagos.}

\noindent\textbf{OE-11 Pagos.}
Dise\~nar e implementar el m\'odulo de liquidaci\'on que permita registrar el estado de pago
de cada comensal (individual y mesa compartida), controlar la concurrencia mediante
\texttt{SELECT FOR UPDATE} en PostgreSQL, soportar la aplicaci\'on de descuentos (porcentuales
o fijos) y el registro de reembolsos con trazabilidad al pago original. La integraci\'on con
un agregador de pago v\'ia tarjeta queda fuera del alcance de la versi\'on~1.0.
\textit{Evidencia: Liquidaci\'on exitosa con aplicaci\'on de descuentos y reembolsos verificada
en la base de datos. Comprobante digital generado al completarse todos los pagos.}

\noindent\textbf{Módulo 7 — Analytics y Panel Administrativo}

\noindent\textbf{OE-12 Analytics y Dashboard.}
Implementar el panel de analytics accesible para Admin y Gerente, sustentado en un
pipeline de procesamiento batch mediante Apache Spark y una arquitectura Medallion
(Bronze, Silver, Gold). El sistema debe transformar los datos transaccionales en
tablas anal\'iticas (\textit{Gold}) para visualizar m\'etricas de operaci\'on
filtradas seg\'un el nivel organizacional: mesas activas, velocidades de venta,
ingresos totales, productos m\'as vendidos y pron\'osticos de demanda.
\textit{Evidencia: Dashboard funcional visualizando m\'etricas provenientes de tablas Gold.
Jobs de Spark documentados y ejecutados exitosamente. Separaci\'on de carga verificada:
las consultas anal\'iticas no afectan el rendimiento de la base transaccional.}

\noindent\textbf{Módulo 8 — Seguridad y Calidad}

\noindent\textbf{OE-13 Seguridad y Auditoría.}
Garantizar comunicación cifrada en tránsito (HTTPS/TLS 1.2+) y aislamiento de datos
mediante políticas RLS en Supabase. Implementar un Registro de Auditoría
(\textit{Audit Log}) inmutable que documente cada creación, modificación o
eliminación de entidades críticas (\textit{snapshots} JSON), y documentar el Aviso
de Privacidad conforme a la LFPDPPP.
\textit{Evidencia: Escaneo de seguridad OWASP ZAP sin vulnerabilidades críticas.
Auditoría de políticas RLS completada. Tabla de \texttt{AuditLog} poblada
correctamente con trazabilidad de IP y usuario actor.}

\subsection{Justificación}

La relevancia de este proyecto radica en proponer una soluci\'on tecnol\'ogica que atienda
la problem\'atica descrita desde un enfoque integral. El sistema omnicanal planteado busca
centralizar en una misma plataforma la gesti\'on de mesas, \'ordenes y pagos, permitiendo que
los comensales participen directamente en el registro de su consumo mediante el acceso a mesas
virtuales por c\'odigos QR \cite{koay_qr_2024, toast_contactless_2026}. A trav\'es de esta din\'amica, cada usuario podr\'a seleccionar
productos, visualizar su consumo individual y decidir si desea pagar \'unicamente lo propio,
contribuir a una cuenta compartida o cubrir consumos adicionales, reduciendo la dependencia
de procesos manuales al cierre del servicio.

Adem\'as del beneficio directo para los clientes, la propuesta aporta ventajas importantes
para los restaurantes. Al contar con una base de datos centralizada y trazabilidad de las
transacciones, el sistema facilita el control de la operaci\'on, mejora la comunicaci\'on entre
comensales, meseros y cocina, y permite disponer de informaci\'on estructurada para la generaci\'on
de reportes y m\'etricas administrativas. Esto representa una oportunidad para fortalecer la
toma de decisiones, optimizar procesos internos y proporcionar capacidades avanzadas
de an\'alisis de comportamiento de consumo e inteligencia de negocio mediante el
procesamiento de grandes vol\'umenes de datos \cite{choo_qr_2023}.

Desde la perspectiva del mercado, la digitalizaci\'on en la industria restaurantera ya no es
opcional. La National Restaurant Association (2024) report\'o que las tecnolog\'ias orientadas
a mejorar velocidad, precisi\'on y conveniencia son relevantes tanto para consumidores como para
operadores \cite{nra_state_2023}, y que el inter\'es de los comensales por ordenar y pagar con mayor autonom\'ia
dentro del restaurante ha crecido de forma sostenida \cite{gonzalez_canirac_2024}.
A nivel global, el mercado de tecnolog\'ia para restaurantes ha registrado un crecimiento
significativo, impulsado por la digitalizaci\'on acelerada del sector tras la pandemia \cite{gonzalez_canirac_2024, nra_state_2023}.

La originalidad del trabajo se encuentra en la integraci\'on de distintas funciones que
usualmente aparecen separadas en soluciones existentes. Mientras algunas plataformas se
especializan en el punto de venta, otras en pedidos digitales y otras en divisi\'on de gastos,
la propuesta desarrollada busca reunir la gesti\'on operativa del restaurante con la interacci\'on
directa del cliente y la flexibilidad en el proceso de pago dentro de un mismo ecosistema
tecnol\'ogico, potenciado por un motor de procesamiento distribuido para la
generaci\'on de inteligencia de negocio avanzada \cite{iovescu_realtime_2024}.

Asimismo, el proyecto presenta un nivel de complejidad significativo, ya que involucra
conocimientos de desarrollo web, bases de datos, sincronizaci\'on en tiempo real, modelado de
procesos, integraci\'on de pasarelas de pago y el despliegue de arquitecturas de datos anal\'iticos con
Apache Spark para el procesamiento y an\'alisis de grandes vol\'umenes de
informaci\'on, constituyendo una aplicaci\'on pr\'actica e integradora de los
conocimientos adquiridos en la formaci\'on acad\'emica.

\subsection{Alcances y Limitaciones}

\subsubsection{Funcionalidades Incluidas en la Versión 1.0}

El presente apartado delimita con precisión el conjunto de funcionalidades, módulos,
usuarios y restricciones técnicas que comprenden el sistema en su versión~1.0.
Todo lo listado constituye un entregable comprometido del proyecto. La
Tabla~\ref{tab:alcances} resume los alcances funcionales confirmados.

{\small
\setlength{\tabcolsep}{4pt}
\renewcommand{\arraystretch}{1.5}
\begin{longtable}{>{\centering\bfseries\arraybackslash}p{1.2cm}
                >{\raggedright\arraybackslash}p{14.2cm}}
\caption{Alcances Funcionales del Sistema v1.0} \label{tab:alcances} \\
\toprule
\thead{ID} & \thead{Funcionalidad incluida} \\
\midrule
\endfirsthead
\caption[]{TABLA \thetable{} (Continuación)} \\
\toprule
\thead{ID} & \thead{Funcionalidad incluida} \\
\midrule
\endhead
\bottomrule
\endfoot
AL-01 & \textbf{Gesti\'on de estructura organizacional jer\'arquica.} Creaci\'on y administraci\'on de los tres niveles: Cadena, Zona y Sucursal, con integridad referencial y herencia de configuraciones (moneda, impuestos, propinas). \\
\rowcolor{altrow}
AL-02 & \textbf{Gobernanza de Datos y Control de Acceso (RBAC).} Matriz de permisos detallada por rol (Admin, Gerente, Mesero, Cocina) con pol\'iticas de seguridad a nivel de fila (RLS) en Supabase. \\
AL-03 & \textbf{Seguridad de Identidad y Contrase\~nas.} Alta de colaboradores con generaci\'on de contrase\~na temporal y flujo obligatorio de cambio en el primer inicio de sesi\'on (\texttt{isFirstLogin}). \\
\rowcolor{altrow}
AL-04 & \textbf{Cat\'alogo de Men\'u Multinivel.} Gesti\'on de \textit{Templates} a nivel de cadena para estandarizaci\'on y men\'us locales por sucursal con categor\'ias, productos, variantes y modificadores. \\
AL-05 & \textbf{Gesti\'on de Mesas y Plano Interactivo (\textit{Floor Map}).} Representaci\'on visual del sal\'on mediante coordenadas $X, Y$ y formas geom\'etricas (\texttt{rect}, \texttt{circle}); visualizaci\'on en tiempo real del estatus ocupacional. \\
\rowcolor{altrow}
AL-06 & \textbf{Mesa Virtual y Seguridad QR.} Generaci\'on de c\'odigos QR con rotaci\'on de tokens de seguridad (\texttt{qrTokenHash}) y versi\'on para invalidaci\'on remota; activaci\'on operativa por el personal. \\
AL-07 & \textbf{Interacci\'on Colaborativa de Comensales.} Flujo de uni\'on de comensales (\textit{Guests}) a una sesi\'on mediante clave de acceso; soporte para \'ordenes simult\'aneas sobre una misma mesa compartida. \\
\rowcolor{altrow}
AL-08 & \textbf{Gesti\'on de \'Ordenes con Integridad Hist\'orica.} Captura de \'ordenes individuales y compartidas con generaci\'on de \textit{snapshots} inmutables de nombre y precio para evitar inconsistencias por cambios futuros en el cat\'alogo. \\
AL-09 & \textbf{Personalizaci\'on de Productos y Modificadores.} Soporte para grupos de modificadores (\textit{Extras}, \textit{T\'erminos}) con reglas de selecci\'on m\'inima/m\'axima y ajustes de precio autom\'aticos. \\
\rowcolor{altrow}
AL-10 & \textbf{Cocina Digital e Inteligencia Operativa (KDS).} Enrutamiento de comandas por estaciones de preparaci\'on; tablero en tiempo real con historial de transiciones de estado para m\'etricas de desempe\~no. \\
AL-11 & \textbf{Divisi\'on de Cuenta Flexible (\textit{Split Bill}).} Liquidaci\'on por art\'iculos seleccionados (\textit{By Item}), partes iguales (\textit{Equal}) o cuenta completa (\textit{Full}), garantizando la consistencia financiera de la sesi\'on. \\
\rowcolor{altrow}
AL-12 & \textbf{M\'odulo Financiero: Pagos, Reembolsos y Descuentos.} Gesti\'on transaccional de pagos multimoneda, aplicaci\'on de descuentos autorizados (fijos o porcentuales) y registro de reembolsos con trazabilidad al proveedor. \\
AL-13 & \textbf{Control de Concurrencia Transaccional.} Mecanismo \texttt{SELECT FOR UPDATE} para prevenir cobros duplicados en pagos simult\'aneos y asegurar la liberaci\'on correcta de la mesa (\texttt{Liquidada}). \\
\rowcolor{altrow}
AL-14 & \textbf{Trazabilidad y Auditor\'ia de Sistemas.} Registro inmutable (\textit{Audit Log}) de cada operaci\'on administrativa, incluyendo direcci\'on IP, agente de usuario y \textit{snapshots} del estado previo y posterior al cambio. \\
AL-15 & \textbf{Pipeline Anal\'itico Apache Spark (Arquitectura Medallion).} Jobs batch para la ingesta y transformaci\'on de datos en capas Bronze, Silver y Gold para an\'alisis de KPI estrat\'egicos. \\
\rowcolor{altrow}
AL-16 & \textbf{Dashboard Administrativo de Inteligencia de Negocio.} Visualizaci\'on de m\'etricas avanzadas (velocidad de venta, ingresos por zona, saturaci\'on de estaciones) extra\'idas del entorno OLAP en el VPS y despliegue en Vercel + Supabase con ambientes Prod y Preview separados. \\
\end{longtable}
}

\newpage
\subsubsection{Límites del Sistema — Funcionalidades Excluidas}

La Tabla~\ref{tab:limites} describe explícitamente las funcionalidades que quedan
fuera del alcance de la versión~1.0, con su justificación correspondiente.

{\small
\setlength{\tabcolsep}{4pt}
\renewcommand{\arraystretch}{1.5}
\begin{longtable}{>{\centering\bfseries\arraybackslash}p{1.2cm}
                >{\raggedright\arraybackslash}p{5.5cm}
                >{\raggedright\arraybackslash}p{8.5cm}}
\caption{Límites del Sistema v1.0 — Funcionalidades Excluidas} \label{tab:limites} \\
\toprule
\thead{ID} & \thead{Funcionalidad excluida} & \thead{Justificación} \\
\midrule
\endfirsthead
\caption[]{TABLA \thetable{} (Continuación)} \\
\toprule
\thead{ID} & \thead{Funcionalidad excluida} & \thead{Justificación} \\
\midrule
\endhead
\bottomrule
\endfoot
LM-01 & Múltiples pisos o áreas por sucursal &
  Requiere modelo de datos y plano de planta más complejo. Pospuesto para versión futura. \\
\rowcolor{altrow}
LM-02 & Recuperación de contraseña de usuarios colaboradores &
  El restablecimiento se gestiona manualmente por el usuario con permisos de alta. Se incluirá en v1.1. \\
LM-03 & Motor de recomendaciones complejo &
  Requiere integraci\'on de modelos de grafos o filtrado colaborativo avanzado. Pospuesto para versiones posteriores. \\
\rowcolor{altrow}
LM-05 & Gestión de inventario en cocina &
  La disponibilidad de productos se gestiona manualmente activando o desactivando ítems en el catálogo. \\
LM-06 & Reservaciones y pre-orden de mesas &
  La activación de la mesa virtual ocurre únicamente cuando el comensal está físicamente en el restaurante. \\
\rowcolor{altrow}
LM-07 & Aplicación móvil nativa (iOS / Android) &
  El sistema se desarrolla como PWA accesible desde el navegador. Sin apps en App Store ni Google Play. \\
LM-08 & Integración con sistemas POS externos &
  La plataforma opera de forma independiente. Sin integración con Micros, Toast, Square, etc. \\
\rowcolor{altrow}
LM-09 & Notificaciones push nativas fuera de la app &
  Las notificaciones se implementan vía Supabase Realtime dentro de la app abierta. Sin push del SO. \\
LM-10 & Programa de lealtad o fidelización de comensales &
  Los datos de consumo se registran para fines analíticos internos, no para programas al cliente final. \\
\rowcolor{altrow}
LM-11 & Agregador de pago v\'ia tarjeta (PCI~DSS) &
  Procesar cobros con tarjeta sin certificaci\'on PCI~DSS expone a la empresa a sanciones de redes de pago y responsabilidad civil. Su integraci\'on en versiones futuras requerir\'a proveedor Level~1 PCI~DSS o auto-certificaci\'on. \\
\end{longtable}
}

\subsubsection{Restricciones Técnicas Asumidas}

La Tabla~\ref{tab:restricciones_asumidas} documenta las restricciones técnicas
que han sido identificadas, evaluadas y asumidas como parte del diseño de la
versión~1.0.

{\small
\setlength{\tabcolsep}{4pt}
\renewcommand{\arraystretch}{1.5}
\begin{longtable}{>{\bfseries\raggedright\arraybackslash}p{4.5cm}
                >{\raggedright\arraybackslash}p{10.5cm}}
\caption{Restricciones Técnicas Asumidas en la Versión 1.0} \label{tab:restricciones_asumidas} \\
\toprule
\thead{Restricción} & \thead{Descripción y justificación} \\
\midrule
\endfirsthead
\caption[]{TABLA \thetable{} (Continuación)} \\
\toprule
\thead{Restricción} & \thead{Descripción y justificación} \\
\midrule
\endhead
\bottomrule
\endfoot
Un solo piso por sucursal &
  El m\'odulo de mesas f\'isicas gestiona una \'unica \'area plana por sucursal. No se implementa un sistema de coordenadas o plano de planta para m\'ultiples pisos o secciones. \\
\rowcolor{altrow}
Sin soporte offline &
  La plataforma requiere conectividad a internet en todo momento. No se implementa sincronizaci\'on offline ni Service Worker de cach\'e de datos transaccionales. \\
Un solo idioma (espa\~nol) &
  La interfaz y todos los textos est\'an en espa\~nol. No se implementa internacionalizaci\'on (i18n) ni soporte multi-idioma. \\
\rowcolor{altrow}
Acceso mediante navegador web &
  Los comensales acceden desde el navegador de su dispositivo m\'ovil. La PWA puede agregarse a la pantalla de inicio, pero no requiere instalaci\'on nativa. \\
Sin agregador de pago v\'ia tarjeta en v1.0 &
  La liquidaci\'on registra el estado de pago por comensal, pero el cobro con tarjeta se gestiona fuera de la plataforma. La integraci\'on de un gateway PCI~DSS queda documentada como trabajo futuro (LM-11). \\
\end{longtable}
}

% ============================================================
%  CAPÍTULO 2 — MARCO TEÓRICO Y ESTADO DEL ARTE
% ============================================================
\newpage
